% Part: lambda-calculus
% Chapter: syntax
% Section: alpha

\documentclass[../../../include/open-logic-section]{subfiles}

\begin{document}

\olfileid{lam}{syn}{alp}
\olsection{$\alpha$-परिवर्तन}

$\lambd[x][x]$ आणि $\lambd[y][y]$ यांचा संबंध काय?
ही दोन्ही पदे ओळख फलन दर्शवतात. अर्थात, विन्यासाच्या
दृष्टीने ती भिन्न पदे आहेत. त्यांच्यात फक्त बद्ध चराच्या
नावाचा फरक आहे; एकातील बद्ध चराचे ``नामांतर'' केल्यावर
दुसरे पद मिळते. याला \emph{$\alpha$-परिवर्तन} म्हणतात.

\begin{defn}[बद्ध चराचे नामांतर, $\aconvone$]
  पद $M$ मध्ये $\lambd[x][N]$ ची आस्थिती असेल,
  $x \neq y$, $y \notin \FV{N}$ असेल आणि
  $\Subst{N}{y}{x}$ परिभाषित असेल, तर ती आस्थिती
  \begin{equation*}
    \lambd[y][\Subst{N}{y}{x}]
  \end{equation*}
  ने बदलून मिळणाऱ्या $M'$ ला \emph{बद्ध चराचे नामांतर}
  म्हणतात आणि $M \redone[\alpha] M'$ असे लिहितात.
\end{defn}

\begin{defn}[संबंधाची रचनांशी सुसंगती]
  पदांवरील संबंध $R$ पुढील अटी पूर्ण करत असेल, तर तो
  \emph{रचनांशी सुसंगत} म्हणतात:
  \begin{enumerate}
  \item $R N N'$ असेल, तर $R \lambd[x][N] \lambd[x][N']$.
  \item $R P P'$ असेल, तर $R (PQ) (P'Q)$.
  \item $R Q Q'$ असेल, तर $R (PQ) (PQ')$.
  \end{enumerate}
\end{defn}

आता ही व्याख्या दुसऱ्या रीतीने मांडू:
\begin{defn}[बद्ध चराचे नामांतर, $\aconvone$]
  \emph{बद्ध चराचे नामांतर} ($\redone[\alpha]$) हा पदांवरील
  सर्वांत लहान रचनांशी सुसंगत संबंध आहे जो पुढील अट पूर्ण करतो:
  \begin{align*}
    &\lambd[x][N] \redone[\alpha] \lambd[y][\Subst{N}{y}{x}] && \text{जर
      $x \neq y$, $y \notin \FV{N}$} \\
    & &&\text{आणि $\Subst{N}{y}{x}$ परिभाषित असेल}
  \end{align*}
\end{defn}

येथे ``सर्वांत लहान'' म्हणजे सुसंगती आणि अतिरिक्त अट
यांमुळे आवश्यक ठरणाऱ्या जोड्याच या संबंधात असतात;
इतर कोणत्याही नसतात. म्हणून हा संबंध पुढीलप्रमाणेही
परिभाषित करता येतो:

\begin{defn}[बद्ध चराचे नामांतर, $\aconvone$] \ollabel{defn:aconvone}
  \emph{बद्ध चराचे नामांतर} ($\aconvone$) आगमनाने
  पुढीलप्रमाणे परिभाषित करतात:
  \begin{enumerate}
  \item $N \aconvone N'$ असेल, तर $\lambd[x][N] \aconvone
    \lambd[x][N']$. \ollabel{defn:aconvone1}
  \item $P \aconvone P'$ असेल, तर $(PQ) \aconvone (P'Q)$. \ollabel{defn:aconvone2}
  \item $Q \aconvone Q'$ असेल, तर $(PQ) \aconvone (PQ')$. \ollabel{defn:aconvone3}
  \item $x \neq y$, $y \notin \FV{N}$ आणि $\Subst{N}{y}{x}$
    परिभाषित असेल, तर
    $\lambd[x][N] \redone[\alpha] \lambd[y][\Subst{N}{y}{x}]$.
    \ollabel{defn:aconvone4}
  \end{enumerate}
\end{defn}

या व्याख्या सममूल्य आहेत; त्यांची सिद्धता सरावासाठी ठेवतो.
यापुढे आपण आगमनाधारित व्याख्या वापरू.

\begin{defn}[$\alpha$-परिवर्तन, $\aconv$]
  \emph{$\alpha$-परिवर्तन} ($\aconv$) हा पदांवरील
  $\aconvone$ समाविष्ट करणारा सर्वांत लहान परावर्ती
  आणि संक्रमक संबंध आहे.
\end{defn}

वरच्याप्रमाणे, ``सर्वांत लहान'' म्हणजे संक्रमणशीलता आणि
$\aconvone$ यांमुळे आवश्यक ठरणाऱ्या जोड्याच या संबंधात
असतात. यावरून पुढील सममूल्य व्याख्या मिळते:

\begin{defn}[$\alpha$-परिवर्तन, $\aconv$] \ollabel{defn:aconv}
  \emph{$\alpha$-परिवर्तन} ($\aconv$) आगमनाने
  पुढीलप्रमाणे परिभाषित करतात:
  \begin{enumerate}
  \item $P \aconv Q$ आणि $Q \aconv R$ असतील, तर $P \aconv R$.
    \ollabel{defn:aconv1}
  \item $P \aconvone Q$ असेल, तर $P \aconv Q$. \ollabel{defn:aconv2}
  \item $P \aconv P$. \ollabel{defn:aconv3}
  \end{enumerate}
\end{defn}

\begin{ex}
  $\lambd[x][f x]$ चे $\alpha$-परिवर्तन करून
  $\lambd[y][f
  y]$ मिळते; उलट दिशेनेही परिवर्तन करता येते.
  अनौपचारिकपणे, ही दोन्ही पदे युक्तिवाद स्वीकारून
  त्यावर $f$ लावल्याचा परिणाम देणारी फलने आहेत;
  त्यांतील~$f$ हा परिसरातील चराचा संदर्भ आहे.

  $\lambd[x][f x]$ चे $\alpha$-परिवर्तन करून
  $\lambd[x][g
    x]$ मिळत नाही. ही पदे अनुक्रमे परिसरातील
  $f$ आणि~$g$ यांचा संदर्भ घेतात आणि म्हणून भिन्न
  फलने दर्शवतात: $f$ आणि $g$ भिन्न असलेल्या परिसरांत
  त्यांचे वर्तन भिन्न असते.
\end{ex}

\begin{prob}
  पुढील पदजोड्या $\alpha$-परिवर्तनीय आहेत का?
  \begin{enumerate}
  \item $\lambd[x][\lambd[y][x]]$ आणि $\lambd[y][\lambd[x][y]]$
  \item $\lambd[x][\lambd[y][x]]$ आणि $\lambd[c][\lambd[b][a]]$
  \item $\lambd[x][\lambd[y][x]]$ आणि $\lambd[c][\lambd[b][a]]$
  \end{enumerate}
\end{prob}

\begin{lem}\ollabel{lem:fv-one}
  $P \aconvone Q$ असेल, तर $\FV{P} = \FV{Q}$.
\end{lem}

\begin{proof}
  $P \aconvone Q$ च्या !!{derivation} वर विगमन करू.
  \begin{enumerate}
  \item शेवटचा नियम \olref{defn:aconvone4} असेल, तर $P$
    हे $\lambd[x][N]$ या रूपाचे आणि $Q$ हे
    $\lambd[y][\Subst{N}{y}{x}]$ या रूपाचे असते;
    येथे $x \neq y$, $y \notin \FV{N}$ आणि
    $\Subst{N}{y}{x}$ परिभाषित असते. $x \in \FV{N}$
    आहे की नाही यानुसार प्रकरणे वेगळी करू:
    \begin{enumerate}
    \item $x \in \FV{N}$ असेल, तर:
      \begin{align*}
        \FV{\lambd[y][\Subst{N}{y}{x}]} &=
        \FV{\Subst{N}{y}{x}} \setminus \{y\} \\
        & = ((\FV{N} \setminus \{x\}) \cup \{y\}) \setminus \{y\}
         && \text{\olref[sub]{thm:infv} नुसार} \\
        & = \FV{N} \setminus \{x\} \\
        & = \FV{\lambd[x][N]}
      \end{align*}
    \item $x \notin \FV{N}$ असेल, तर:
      \begin{align*}
        \FV{\lambd[y][\Subst{N}{y}{x}]}
        & = \FV{\Subst{N}{y}{x}} \setminus \{y\} \\
        & = \FV{N} \setminus \{x\}
         && \text{\olref[sub]{thm:notinfv} नुसार} \\
        & = \FV{\lambd[x][N]}.
      \end{align*}
    \end{enumerate}
  \item उरलेली तीन प्रकरणे सरावासाठी आहेत.
  \end{enumerate}
\end{proof}

\begin{prob}
  \olref[lam][syn][alp]{lem:fv-one} ची सिद्धता पूर्ण करा.
\end{prob}

\begin{lem}\ollabel{lem:inv}
  $P \aconvone Q$ असेल, तर $Q \aconvone P$.
\end{lem}

\begin{proof}
  $P \aconvone Q$ च्या !!{derivation} वर विगमन करू.
  \begin{enumerate}
  \item शेवटचा नियम \olref{defn:aconvone4} असेल, तर $P$
    हे $\lambd[x][N]$ या रूपाचे आणि $Q$ हे
    $\lambd[y][\Subst{N}{y}{x}]$ या रूपाचे असते;
    येथे $x \neq y$, $y \notin \FV{N}$ आणि
    $\Subst{N}{y}{x}$ परिभाषित असते. प्रथम,
    \olref[sub]{thm:clr} नुसार $x \notin
    \FV{\Subst{N}{y}{x}}$ आहे. \olref[sub]{thm:inv}
    नुसार $\Subst{\Subst{N}{y}{x}}{x}{y}$ केवळ
    परिभाषितच नाही, तर~$N$ च्या बरोबरही आहे. म्हणून
    \olref{defn:aconvone4} नुसार
    $\lambd[y][\Subst{N}{y}{x}]
    \aconvone \lambd[x][\Subst{\Subst{N}{y}{x}}{x}{y}] =
    \lambd[x][N]$.
  \end{enumerate}
\end{proof}

\begin{prob}
  \olref[lam][syn][alp]{lem:inv} ची सिद्धता पूर्ण करा.
\end{prob}

\begin{thm}
  $\alpha$-परिवर्तन हा पदांवरील तुल्यता संबंध आहे;
  म्हणजे तो परावर्ती, सममित आणि संक्रमक आहे.
\end{thm}

\begin{proof}
  \begin{enumerate}
  \item प्रत्येक पद $M$ साठी, बद्ध चरांच्या \emph{शून्य}
    नामांतरांनी $M$ चे $M$ मध्ये परिवर्तन होते.
  \item $P$ चे बद्ध चरांच्या काही नामांतरांनी $Q$ मध्ये
    $\alpha$-परिवर्तन होत असेल, तर \olref{lem:inv}
    वापरून ती नामांतरे उलट क्रमाने उलटवता येतात आणि
    $Q$ पासून~$P$ मिळते.
  \item $P$ चे बद्ध चरांच्या एका क्रमाने $Q$ मध्ये आणि
    $Q$ चे दुसऱ्या क्रमाने $R$ मध्ये $\alpha$-परिवर्तन
    होत असेल, तर दोन्ही क्रम सलग लावून $P$ पासून
    $R$ मिळते.
  \end{enumerate}
\end{proof}

यापुढे $M$ आणि $N$ ही पदे \emph{$\alpha$-सममूल्य},
$M \aeq N$, आहेत असे म्हणू, जर $M$ चे $N$ मध्ये
$\alpha$-परिवर्तन होत असेल. आपण नुकतेच दाखवल्याप्रमाणे,
हे घडते तेव्हाच आणि तेव्हाच $N$ चेही~$M$ मध्ये
$\alpha$-परिवर्तन होते.

\begin{thm}\ollabel{thm:fv}
  $M \aeq N$ असेल, तर $\FV{M} = \FV{N}$.
\end{thm}

\begin{proof}
  \olref{lem:fv-one} वरून लगेच मिळते.
\end{proof}

\begin{lem}\ollabel{lem:sub:R}
  $R \aeq R'$ असेल आणि $\Subst{M}{R}{y}$ परिभाषित
  असेल, तर $\Subst{M}{R'}{y}$ परिभाषित असते आणि
  $\Subst{M}{R}{y}$ शी $\alpha$-सममूल्य असते.
\end{lem}

\begin{proof}
  सरावासाठी.
\end{proof}

\begin{prob}
  \olref[lam][syn][alp]{lem:sub:R} सिद्ध करा.
\end{prob}

\olref[sub]{sec} मध्ये पाहिल्याप्रमाणे, काही प्रसंगी
आदेशन अपरिभाषित असते. पण पदांवर $\alpha$-परिवर्तन
वापरून बद्ध चरांची नावे बदलल्यास आदेशन नेहमी परिभाषित
करता येते. परिणामी पदाची $\alpha$-सममूल्यता टिकते,
हे पुढील प्रमेय दाखवते:

\begin{thm}\ollabel{thm:sub}
  कोणत्याही $M$, $R$ आणि~$y$ साठी असे $M'$ असते की
  $M \aeq M'$ आणि $\Subst{M'}{R}{y}$ परिभाषित असते.
  आणखी एखादी जोडी $M'' \aeq M$ आणि $R''$ अशी असेल
  की $\Subst{M''}{R''}{y}$ परिभाषित आहे आणि $R'' \aeq R$,
  तर $\Subst{M'}{R}{y} \aeq \Subst{M''}{R''}{y}$.
\end{thm}

\begin{proof}
  $M$ च्या रचनेवर विगमन करू:
  \begin{enumerate}
  \item $M$ हे चर~$z$ आहे: सरावासाठी.
  \item $M$ हे $\lambd[x][N]$ या रूपाचे आहे असे समजू.
    $x$ आणि $y$ यांच्याहून भिन्न, तसेच $z \notin \FV{N}$
    आणि $z
    \notin \FV{R}$ अशी अट पूर्ण करणारा चर $z$ निवडा.
    विगमन गृहीतकानुसार असे $N'$ आहे की $N' \aeq N$
    आणि $\Subst{N'}{z}{x}$ परिभाषित आहे. मग
    \olref{defn:aconvone}\olref{defn:aconvone1} नुसार
    $\lambd[x][N] \aeq \lambd[x][N']$ देखील आहे.
    आता \olref{defn:aconvone}\olref{defn:aconvone4} नुसार
    $\lambd[x][N']
    \aeq \lambd[z][\Subst{N'}{z}{x}]$.
    हे शक्य आहे, कारण $z \ne x$, $z \notin \FV{N'}$
    आणि $\Subst{N'}{z}{x}$ परिभाषित आहे. शेवटी,
    $\Subst{\lambd[z][\Subst{N'}{z}{x}]}{R}{y}$
    परिभाषित आहे, कारण $z \neq y$ आणि $z \notin \FV{R}$.

    त्याच अटी पूर्ण करणारे दुसरे $N''$ आणि $R''$ असल्यास,
    \begin{multline*}
      \Subst{(\lambd[z][\Subst{N''}{z}{x}])}{R''}{y} =\\
    \begin{aligned}
      &= \lambd[z][\Subst{\Subst{N''}{z}{x}}{R''}{y}] \\
      &\aeq \lambd[z][\Subst{\Subst{N''}{z}{x}}{R}{y}] && \text{\olref{lem:sub:R} नुसार}\\
      &\aeq\lambd[z][\Subst{\Subst{N'}{z}{x}}{R}{y}]
      && \text{विगमन गृहीतकानुसार}\\
      &=\Subst{(\lambd[z][\Subst{N'}{z}{x}])}{R}{y}
    \end{aligned}
    \end{multline*}
    \item $M$ हे $(PQ)$ या रूपाचे आहे: सरावासाठी.
  \end{enumerate}
\end{proof}

\begin{prob}
  \olref[lam][syn][alp]{thm:sub} ची सिद्धता पूर्ण करा.
\end{prob}

\begin{cor}\ollabel{cor:sub}
  कोणत्याही $M$, $R$ आणि~$y$ साठी $M'$ आणि $R'$ ची
  अशी जोडी असते की $M' \aeq M$, $R \aeq R'$
  आणि $\Subst{M'}{R'}{y}$ परिभाषित असते. आणखी
  $M'' \aeq M$ आणि $R'' \aeq R$ अशी जोडी असेल
  आणि $\Subst{M''}{R''}{y}$ परिभाषित असेल, तर
  $\Subst{M'}{R'}{y} \aeq
  \Subst{M''}{R''}{y}$.
\end{cor}

\begin{proof}
  \olref{thm:sub} वरून लगेच मिळते.
\end{proof}

\end{document}
